\chapter{Dedekind 分割}
\label{chap:dedekind-construction}

\chapterref{chap:real-number-system}中的每个有理缺口都有一个确定的左部。Dedekind 的办法是把这个左部本身当作新数。这个想法很短，核验却不能省略：包含关系是否给出全序，两个左部怎样相加、相乘，负号与倒数怎样定义，以及一族新数是否仍有上确界，都要从集合定义推出。只有这些步骤完成以后，分割才组成实数系，而不只是若干“缺口的名称”。

本章的证明有两类。第一类说明某个公式确实产生分割，即所得集合非空、不等于 \(\Q\)、向下封闭且没有最大元；第二类验证这些公式满足域公理及序公理。加法几乎逐项继承有理数运算，乘法则必须先处理正数，再扩张符号。确界性质放在最后证明，因为它取决于分割以左部表示这一特定选择。

\section{分割的定义、序与有理数嵌入}

\subsection{分割的动机与严格定义}


\begin{definition}
有理数集的子集 \(\alpha\subseteq\Q\) 称为一个 \term{Dedekind 分割}，如果
\begin{enumerate}
 \item \(\alpha\ne\varnothing\) 且 \(\alpha\ne\Q\)；
 \item 若 \(x\in\alpha\) 且 \(y<x\)，则 \(y\in\alpha\)；
 \item \(\alpha\) 没有最大元。
\end{enumerate}
全体分割组成的集合记为 \(\mathbb D\)。
\end{definition}

这里“全体分割”确实是集合：每个分割都是 \(\Q\) 的子集，而
\(\mathbb D\) 可从幂集 \(\mathcal P(\Q)\) 中按上述三项条件分离出来。这个存在性核对只使用附录A中的受限概括，不使用任何分析学完备性。

分割只记录左部；右部是 \(\Q\setminus\alpha\)。向下封闭性保证每个左部元素小于每个右部元素，无最大元则统一处理“边界为有理数”和“边界为无理数”两种情形。

三个条件各自排除一种退化。若允许空集或全体 \(\Q\)，所谓边界便落在所有有理数之外；若没有向下封闭性，集合中保留的就不是边界左侧的全部信息；若允许最大元，则同一个有理边界 \(r\) 会同时产生 \(\{q:q<r\}\) 与 \(\{q:q\le r\}\) 两种表示。规定左部没有最大元，是为了只保留前一种标准表示。

对 \(r\in\Q\)，定义主分割
\[
 r^*=\{q\in\Q:q<r\}.
\]
\chapterref{chap:real-number-system}的集合 \(L=\{q:q<0\text{ 或 }q^2<2\}\) 也是分割，但不等于任何 \(r^*\)。

这两个例子说明分割比有理数增加了什么。主分割的边界是已有的有理数 \(r\)，而 \(L\) 的左右两部在 \(\Q\) 中没有共同边界。构造并不先给缺口命名；它把每个符合条件的左部都作为一个数，随后再由域运算和完备性判断这些对象是否组成所需的实数系。

\subsection{分割的相等关系与序关系}


分割是集合，因此相等按集合相等定义。规定
\[
 \alpha\le\beta\quad\Longleftrightarrow\quad\alpha\subseteq\beta.
\]

\begin{proposition}
包含关系使 \(\mathbb D\) 成为全序集。
\end{proposition}

\begin{proof}
只需证明任意两个分割可比较。若 \(\alpha\nsubseteq\beta\)，取 \(a\in\alpha\setminus\beta\)。对任意 \(b\in\beta\)，不可能有 \(a<b\)，否则由 \(\beta\) 向下封闭得 \(a\in\beta\)。故 \(b\le a\)。又 \(a\in\alpha\)，所以 \(b\in\alpha\)（若 \(b=a\) 会有 \(a\in\beta\)，故实际为 \(b<a\)）。因此 \(\beta\subseteq\alpha\)。
\end{proof}

严格包含对应严格次序。对有理数 \(r,s\)，有
\[
 r<s\quad\Longleftrightarrow\quad r^*\subsetneq s^*.
\]

\begin{proposition}[分割的基本性质]
设 \(\alpha\in\mathbb D\)。
\begin{enumerate}
 \item 若 \(a\in\alpha\)、\(r\notin\alpha\)，则 \(a<r\)。
 \item \(\alpha>0^*\) 当且仅当 \(0\in\alpha\)，也当且仅当 \(\alpha\) 含有正有理数。
 \item 给定 \(a_0\in\alpha\)、\(r_0\notin\alpha\) 与 \(\varepsilon>0\)，存在
 \[
   a\in\alpha,\qquad r\notin\alpha,
   \qquad a_0\le a<r\le r_0,\qquad r-a<\varepsilon.
 \]
\end{enumerate}
\end{proposition}

\begin{proof}
若 \(r\le a\)，则 \(r=a\) 与 \(r\notin\alpha\) 矛盾；若 \(r<a\)，向下封闭性又给出 \(r\in\alpha\)。故必有 \(a<r\)。

若 \(\alpha>0^*\)，可取 \(x\in\alpha\setminus0^*\)，于是 \(x\ge0\)。当 \(x>0\) 时向下封闭性给出 \(0\in\alpha\)，当 \(x=0\) 时结论已经成立。反过来，\(0\in\alpha\) 说明 \(0^*\subsetneq\alpha\)。又因 \(\alpha\) 没有最大元，\(0\in\alpha\) 时必有某个正有理数属于 \(\alpha\)；含正有理数当然也推出 \(0\in\alpha\)。

由第一项，\(a_0<r_0\)。对区间 \([a_0,r_0]\) 逐次二分：若中点属于 \(\alpha\)，以它替换左端点；否则以它替换右端点。第 \(n\) 次以后仍有左端点 \(a_n\in\alpha\)、右端点 \(r_n\notin\alpha\)，并且
 \[
   r_n-a_n=2^{-n}(r_0-a_0).
 \]
比值 \((r_0-a_0)/\varepsilon\) 是正有理数，取整数 \(n\) 使
\(2^n>(r_0-a_0)/\varepsilon\) 即可；这个选择只用整数的初等 Archimedes 性质。于是 \(r_n-a_n<\varepsilon\)。取 \(a=a_n,r=r_n\) 即得结论。
\end{proof}

第三项没有声称缺口在 \(\Q\) 中有边界；它只说左右两侧可以达到任意给定的有理精度。后面构造加法逆元和乘法逆元时，正是这一夹逼性质把集合定义转化为恒等式。

\subsection{有理数的嵌入}


定义 \(\iota:\Q\to\mathbb D\) 为 \(\iota(r)=r^*\)。

\begin{theorem}
\(\iota\) 是保序单射：
\[
 r<s\Longleftrightarrow\iota(r)<\iota(s).
\]
\end{theorem}

\begin{proof}
前面已经证明 \(r<s\) 当且仅当 \(r^*\subsetneq s^*\)。因此 \(\iota\) 保持严格次序。若 \(\iota(r)=\iota(s)\)，则既不可能 \(r<s\)，也不可能 \(s<r\)，故 \(r=s\)。
\end{proof}

\begin{corollary}
嵌入后的有理数在 \(\mathbb D\) 中稠密：若 \(\alpha<\beta\)，则存在 \(q\in\Q\) 使
\[
 \alpha<q^*<\beta.
\]
\end{corollary}

\begin{proof}
取 \(b\in\beta\setminus\alpha\)。由 \(\beta\) 无最大元，选 \(q\in\beta\) 且 \(q>b\)。因为 \(b\notin\alpha\)，每个 \(a\in\alpha\) 都满足 \(a<b<q\)，故 \(\alpha\subset q^*\)。严格性来自 \(b\in q^*\setminus\alpha\)。又 \(q\in\beta\) 且 \(\beta\) 向下封闭，所以 \(q^*\subset\beta\)；因 \(q\in\beta\setminus q^*\)，包含严格。
\end{proof}

因此可把 \(\Q\) 与其像识别。加法和乘法尚未定义；嵌入保持运算的性质将在相应定义之后证明。

\section{加法和零元}

若 \(\alpha,\beta\) 分别表示两个数的全部有理下近似，那么它们的和应由所有 \(a+b\)（\(a\in\alpha,b\in\beta\)）从下方逼近。这一想法直接给出定义，但仍须验证所得集合没有遗漏下方点，也没有产生最大元。


定义
\[
 \alpha+\beta=\{a+b:a\in\alpha,\ b\in\beta\},\qquad
 0_{\mathbb D}=0^*.
\]

\begin{proposition}
\(\alpha+\beta\) 是分割；加法在 \(\mathbb D\) 上满足交换律、结合律，零元为 \(0^*\)。
\end{proposition}

\begin{proof}
取 \(a\in\alpha,b\in\beta\)，则 \(a+b\) 说明非空。取 \(u\notin\alpha,v\notin\beta\)，则每个 \(a<u,b<v\)，故 \(a+b<u+v\)，所以 \(u+v\) 不属于该集合，证得它不等于 \(\Q\)。若 \(x=a+b\) 且 \(y<x\)，令 \(a'=a+(y-x)\)，则 \(a'<a\)、\(a'\in\alpha\)，且 \(y=a'+b\)。若 \(a+b\) 已给定，取 \(a'>a\) 且 \(a'\in\alpha\)，便有 \(a'+b>a+b\)，故无最大元。

交换律、结合律由有理数相应法则逐元素得到。对任意 \(a\in\alpha\)、\(q<0\)，有 \(a+q<a\)，故 \(\alpha+0^*\subseteq\alpha\)。反之，对 \(x\in\alpha\)，取 \(a\in\alpha\) 且 \(a>x\)，则 \(x=a+(x-a)\) 且 \(x-a<0\)，故 \(x\in\alpha+0^*\)。
\end{proof}

\begin{proposition}
对任意 \(r,s\in\Q\)，有 \(r^*+s^*=(r+s)^*\)。
\end{proposition}

\begin{proof}
若 \(x=a+b\)，其中 \(a<r,b<s\)，则 \(x<r+s\)。反之，若 \(x<r+s\)，取有理数 \(a\) 满足 \(x-s<a<r\)，令 \(b=x-a\)，便有 \(a<r\)、\(b<s\) 且 \(x=a+b\)。
\end{proof}

\section{负元与加法群}

把 \(\alpha\) 中每个有理数直接取相反数会得到向上封闭的集合，不能作为分割。负元的左部应由原分割右部的相反数从下方生成；定义中的严格不等式同时保证向下封闭和无最大元。

定义加法逆元
\[
 -\alpha=\{q\in\Q:\text{存在 }r\notin\alpha\text{ 使 }q<-r\}.
\]

\begin{theorem}
\(-\alpha\) 是分割，且 \(\alpha+(-\alpha)=0^*\)。
\end{theorem}

\begin{proof}
取 \(r_0\notin\alpha\)，任意 \(q<-r_0\) 都属于 \(-\alpha\)，故该集合非空。再固定 \(a_0\in\alpha\)。若 \(q\in-\alpha\)，按定义存在 \(r\notin\alpha\) 使 \(q<-r\)；分割的基本性质给出 \(a_0<r\)，所以 \(q<-a_0\)。因此 \(-a_0\notin-\alpha\)，该集合不等于 \(\Q\)。

若 \(q\in-\alpha\)、\(t<q\)，原来证明 \(q\in-\alpha\) 的同一个右部元素也证明 \(t\in-\alpha\)，故集合向下封闭。若 \(q\in-\alpha\)，取 \(r\notin\alpha\) 使 \(q<-r\)，再由有理数稠密性选取 \(t\) 满足 \(q<t<-r\)，则 \(t\in-\alpha\)。所以 \(-\alpha\) 没有最大元，确为分割。

若 \(a\in\alpha\)、\(q\in-\alpha\)，取 \(r\notin\alpha\) 使 \(q<-r\)。由 \(a<r\)，得 \(a+q<a-r<0\)，所以 \(\alpha+(-\alpha)\subseteq0^*\)。

反之，取 \(x<0\)，并任取 \(a_0\in\alpha\)、\(r_0\notin\alpha\)。对精度 \(-x>0\) 使用前面的分割夹逼命题，得到 \(a\in\alpha,r\notin\alpha\) 且 \(r-a<-x\)。令 \(q=x-a\)，则
 \[
   q=x-a<-r,
 \]
所以 \(q\in-\alpha\)，并且 \(x=a+q\in\alpha+(-\alpha)\)。这对每个 \(x<0\) 成立，故 \(0^*\subseteq\alpha+(-\alpha)\)。
\end{proof}

于是 \((\mathbb D,+)\) 是交换群，而且 \(-(r^*)=(-r)^*\)。有理数嵌入保持加法和相反数。

\section{正分割的乘法}

乘法比加法多一个符号问题。只有在两个数均为正时，“下近似的乘积仍是下近似”才保持次序，因此先在正分割上定义乘法，再用符号法则扩张。非正有理数自动放入积的左部，是因为正数的每个非正数都位于其左侧。


称 \(\alpha>0^*\) 为正分割，记 \(\alpha_+=\alpha\cap(0,\infty)\)。对正分割定义
\[
 \alpha\beta=
 \{q\in\Q:q\le0\text{，或存在 }a\in\alpha_+,b\in\beta_+
 \text{ 使 }q<ab\}.
\]
\begin{lemma}[正分割乘法演算]
若 \(\alpha,\beta,\gamma>0^*\)，则 \(\alpha\beta\) 是正分割，并且
\[
 \alpha\beta=\beta\alpha,\qquad
 (\alpha\beta)\gamma=\alpha(\beta\gamma),\qquad
 \alpha 1^*=\alpha,\qquad
 \alpha(\beta+\gamma)=\alpha\beta+\alpha\gamma.
\]
\end{lemma}

\begin{proof}
积包含所有非正有理数，故非空。取 \(u\notin\alpha\)、\(v\notin\beta\)。由于两个分割均为正，\(u,v>0\)，且每个 \(a\in\alpha_+\)、\(b\in\beta_+\) 都满足 \(a<u\)、\(b<v\)。于是 \(ab<uv\)，从而 \(uv\notin\alpha\beta\)，积不是全体 \(\Q\)。

若 \(q\in\alpha\beta\) 且 \(t<q\)，则 \(q\le0\) 时 \(t\le0\)；若 \(q>0\)，存在 \(a\in\alpha_+,b\in\beta_+\) 使 \(q<ab\)，于是同样有 \(t<ab\)。故积向下封闭。再证无最大元。若 \(q<0\)，则 \(0\in\alpha\beta\) 且 \(q<0\)；若 \(q=0\)，任选正的 \(a\in\alpha,b\in\beta\)，在 \(0\) 与 \(ab\) 之间取有理数；若 \(q>0\)，先取 \(a,b\) 使 \(q<ab\)，再取 \(a'\in\alpha_+\) 且 \(a'>a\)，便可在 \(q\) 与 \(a'b\) 之间取更大的积内元素。因此 \(\alpha\beta\) 是正分割。

交换律来自 \(ab=ba\)。以下逐项证明结合律；非正有理数本来属于两边，只需比较正有理数。若 \(q\in(\alpha\beta)\gamma\) 且 \(q>0\)，则存在 \(u\in(\alpha\beta)_+\)、\(c\in\gamma_+\) 使 \(q<uc\)。再由 \(u\in\alpha\beta\) 取 \(a\in\alpha_+,b\in\beta_+\) 使 \(u<ab\)，于是 \(q<abc\)。在 \(q/a\) 与 \(bc\) 之间取有理数 \(v\)，则 \(v\in\beta\gamma\) 且 \(q<av\)，故 \(q\in\alpha(\beta\gamma)\)。

反过来，若正数 \(q\in\alpha(\beta\gamma)\)，可取 \(a\in\alpha_+\)、\(v\in(\beta\gamma)_+\) 使 \(q<av\)，再取 \(b\in\beta_+,c\in\gamma_+\) 使 \(v<bc\)。于是 \(q<abc\)。在 \(q/c\) 与 \(ab\) 之间取有理数 \(u\)，便有 \(u\in\alpha\beta\) 且 \(q<uc\)，从而 \(q\in(\alpha\beta)\gamma\)。两个包含方向均已得到。

再证单位元。若正数 \(q\in\alpha1^*\)，则 \(q<ab\)，其中 \(a\in\alpha_+\)、\(0<b<1\)，所以 \(q<a\) 且 \(q\in\alpha\)。反之，若 \(q\in\alpha\) 且 \(q>0\)，取 \(a\in\alpha\) 使 \(a>q\)，再取有理数 \(b\) 使 \(q/a<b<1\)，便有 \(q<ab\)，所以 \(q\in\alpha1^*\)。非正部分相同，故 \(\alpha1^*=\alpha\)。

最后证明分配律。若正数 \(q\in\alpha(\beta+\gamma)\)，取 \(a\in\alpha_+\) 与 \(v=b+c\in(\beta+\gamma)_+\) 使 \(q<av\)，其中 \(b\in\beta,c\in\gamma\)。若 \(ab\le0\)，则 \(ab\) 按定义属于 \(\alpha\beta\)；若 \(ab>0\)，取 \(a'\in\alpha\) 且 \(a'>a\)，便有 \(ab<a'b\)，所以 \(ab\in\alpha\beta\)。同理 \(ac\in\alpha\gamma\)。因此
 \[
   q<a(b+c)=ab+ac\in\alpha\beta+\alpha\gamma,
 \]
而右边是向下封闭的分割，故 \(q\) 属于右边。

反过来，非正的 \(q\) 已自动属于左边，设 \(q>0\) 且 \(q=x+y\)，其中 \(x\in\alpha\beta\)、\(y\in\alpha\gamma\)。可分别选正数 \(a,a'\in\alpha\)、\(b\in\beta\)、\(c\in\gamma\)，使 \(x<ab\)、\(y<a'c\)；当 \(x\le0\) 或 \(y\le0\) 时，任取相应的正因子也能得到严格不等式。再取 \(d\in\alpha_+\) 且 \(d>\max\{a,a'\}\)，便有
\[
  q=x+y<ab+a'c<d(b+c).
\]
因 \(b+c\in\beta+\gamma\)，得到 \(q\in\alpha(\beta+\gamma)\)。这证明两边相等。
\end{proof}

\begin{proposition}
若 \(r,s>0\) 为有理数，则 \(r^*s^*=(rs)^*\)。
\end{proposition}

\begin{proof}
若 \(q\in r^*s^*\) 且 \(q>0\)，则存在 \(a<r,b<s\) 使 \(q<ab<rs\)，所以 \(q\in(rs)^*\)。反之，若 \(0<q<rs\)，先取正有理数 \(a<r\) 使 \(q/a<s\)，再取 \(b\) 满足 \(q/a<b<s\)。于是 \(q<ab\)。非正有理数按定义属于两个分割。
\end{proof}

\section{乘法的符号扩张}

包含次序与加法已经足以判断相反数的符号。若 \(\alpha<\beta\)，则
\(\alpha+\gamma<\beta+\gamma\)：非严格包含由加法定义得到；若平移后相等，再加上 \(-\gamma\) 就会推出 \(\alpha=\beta\)。特别地，
\[
  \alpha<0^*\quad\Longleftrightarrow\quad-\alpha>0^*.
\]
因此，每个非零分割都能唯一写成一个正分割或正分割的相反数。

若一因子为 \(0^*\)，定义其积为 \(0^*\)。对非零分割置
\[
 |\alpha|=\begin{cases}
 \alpha,&\alpha>0^*,\\
 -\alpha,&\alpha<0^*,
 \end{cases}
\]
再规定同号相乘取 \(|\alpha||\beta|\)，异号相乘取
\(-(|\alpha||\beta|)\)。

\begin{proposition}
上述公式对任意 \(\alpha,\beta\in\mathbb D\) 定义了一个分割，并且乘法满足交换律、结合律，以 \(1^*\) 为单位元，以 \(0^*\) 为吸收元。
\end{proposition}

\begin{proof}
非零分割的绝对值为正分割，故 \(|\alpha||\beta|\) 已由上一节定义并且仍是正分割；它的相反数也是分割。因此四种符号组合都给出 \(\mathbb D\) 中的唯一元素，含零的情形则按定义处理。

交换律由正分割乘法的交换律和“同号或异号”这一分类的对称性得到。验证结合律时，只要有一项为零，两边都为 \(0^*\)；三项均非零时，两边的绝对值都是
\[
  (|\alpha||\beta|)|\gamma|
  =|\alpha|(|\beta||\gamma|),
\]
而两边的符号都由三个因子中负号的个数是偶数还是奇数决定。因此两边相等。因 \(1^*>0^*\)，正分割的单位元公式给出 \(\alpha1^*=\alpha\)：\(\alpha>0^*\) 时直接使用上一节的结论，\(\alpha<0^*\) 时两边同时带负号，\(\alpha=0^*\) 时按定义成立。吸收元也已经写入定义。
\end{proof}

主分割的正数乘法已在上一节核对。再结合零与符号两种情形，得到
\[
 r^*s^*=(rs)^*\qquad(r,s\in\Q).
\]
所以有理数嵌入不仅保序、保持加法，也保持乘法。

\section{分配律}

\begin{proposition}[一般分配律]
对任意分割 \(\alpha,\beta,\gamma\)，有
\[
  \alpha(\beta+\gamma)=\alpha\beta+\alpha\gamma.
\]
\end{proposition}

\begin{proof}
正分割乘法演算已经证明三项均为正的情形；若其中某项为零，结论直接成立。先固定 \(\alpha>0^*\)。若 \(\beta,\gamma<0^*\)，写 \(\beta=-\beta_0\)、\(\gamma=-\gamma_0\)，其中 \(\beta_0,\gamma_0>0^*\)，则由正数情形和加法逆元唯一性，
\[
 \alpha(\beta+\gamma)
 =-\alpha(\beta_0+\gamma_0)
 =-(\alpha\beta_0+\alpha\gamma_0)
 =\alpha\beta+\alpha\gamma.
\]

若符号相反，不妨 \(\beta>0^*\)、\(\gamma=-\delta<0^*\)。当 \(\beta\ge\delta\) 时令 \(\eta=\beta-\delta\ge0^*\)。正数情形给出
\[
  \alpha\beta=\alpha(\eta+\delta)=\alpha\eta+\alpha\delta.
\]
在加法群中消去 \(\alpha\delta\)，得到
\[
  \alpha(\beta-\delta)=\alpha\beta-\alpha\delta.
\]
这正是所需公式。若 \(\beta<\delta\)，对正数 \(\delta-\beta\) 应用刚证的等式，再取相反数即可。

至此第一因子为正时均已证明。若 \(\alpha<0^*\)，写 \(\alpha=-\alpha_0\)，先对 \(\alpha_0\) 使用已证公式，再整体取相反数。故分配律对任意分割成立。
\end{proof}

\section{乘法逆元}

倒数会反转正数的次序。若只对 \(a\in\alpha_+\) 取 \(1/a\)，得到的是从上方逼近倒数的集合；因此定义必须使用 \(\alpha\) 的右部。对每个右部元素 \(r\)，数 \(1/r\) 位于目标倒数的左侧，再对这些数向下封闭便得到下式。

对正分割 \(\alpha\)，定义
\[
 \alpha^{-1}=
 \{q\in\Q:q\le0\text{，或存在 }r>0,\ r\notin\alpha,\ q<r^{-1}\}.
\]

\begin{theorem}
若 \(\alpha>0^*\)，则 \(\alpha^{-1}\) 是正分割且
\[
 \alpha\alpha^{-1}=1^*.
\]
\end{theorem}

\begin{proof}
先验证 \(\alpha^{-1}\) 是分割。它包含所有非正有理数，因而非空。由于 \(\alpha>0^*\) 且 \(\alpha\ne\Q\)，可取 \(a\in\alpha_+\)。若 \(a^{-1}\in\alpha^{-1}\)，则存在 \(r>0\) 满足 \(r\notin\alpha\) 且 \(a^{-1}<r^{-1}\)，从而 \(r<a\)。这与 \(a\in\alpha\) 及 \(\alpha\) 向下封闭矛盾。因此 \(a^{-1}\notin\alpha^{-1}\)，该集合不是 \(\Q\)。

若 \(q\in\alpha^{-1}\) 且 \(t<q\)，则 \(q\le 0\) 时有 \(t\le0\)；\(q>0\) 时，可取 \(r\notin\alpha\) 使 \(q<r^{-1}\)，于是 \(t<r^{-1}\)。两种情形都给出 \(t\in\alpha^{-1}\)，故它向下封闭。最后说明它没有最大元。对 \(q\le0\)，先取任意正数 \(r\notin\alpha\)，再取 \(0<t<r^{-1}\)，便有 \(q<t\in\alpha^{-1}\)；对 \(q>0\)，由定义取 \(r\notin\alpha\) 使 \(q<r^{-1}\)，再在二者之间取有理数 \(t\)。因此 \(\alpha^{-1}\) 确为正分割。

先证 \(\alpha\alpha^{-1}\subseteq1^*\)。非正有理数已经属于 \(1^*\)。若 \(z>0\) 且 \(z\in\alpha\alpha^{-1}\)，可取 \(a\in\alpha_+\)、\(q\in(\alpha^{-1})_+\) 使 \(z<aq\)。再由 \(q\in\alpha^{-1}\) 取 \(r>0\)、\(r\notin\alpha\) 使 \(q<r^{-1}\)。分割的基本性质给出 \(a<r\)，故
\[
 z<aq<ar^{-1}<1.
\]
于是 \(z\in1^*\)，第一个包含得证。

反之，给定 \(0<x<1\)，取 \(a_0\in\alpha_+\)、\(r_0\notin\alpha\)。对正数
\[
 \varepsilon=\frac{1-x}{x}a_0
\]
使用分割夹逼命题，取得 \(a\in\alpha_+\)、\(r\notin\alpha\)，使
\[
 r-a<\frac{1-x}{x}a_0\le\frac{1-x}{x}a.
\]
整理右端不等式得 \(xr<a\)，即 \(x/a<1/r\)。由有理数稠密性取 \(q\) 使
\[
  \frac{x}{a}<q<\frac1r.
\]
则 \(q\in\alpha^{-1}\) 且 \(x<aq\)，由积的定义得 \(x\in\alpha\alpha^{-1}\)。负数和零本来就在积中，故 \(1^*\subseteq\alpha\alpha^{-1}\)。
\end{proof}

负分割的逆元由
\[
  \alpha^{-1}=-\bigl((-\alpha)^{-1}\bigr)
\]
定义。此时 \(\alpha\) 与 \(\alpha^{-1}\) 同为负分割，按符号扩张的乘法，
\[
  \alpha\alpha^{-1}
  =(-\alpha)(-\alpha)^{-1}
  =1^*.
\]
因此每个非零分割都有乘法逆元。

\section{顺序与域运算的相容性}

\begin{proposition}[序与运算的相容性细化]
对任意分割 \(\alpha,\beta,\gamma\)，有
\[
 \alpha\subseteq\beta
 \Longrightarrow
 \alpha+\gamma\subseteq\beta+\gamma.
\]
若 \(0^*\subseteq\gamma\)，则
\[
 \alpha\subseteq\beta
 \Longrightarrow
 \alpha\gamma\subseteq\beta\gamma.
\]
加法在严格包含下仍严格保序；乘法只有在 \(0^*\subsetneq\gamma\) 时严格保序。若 \(\gamma=0^*\)，两边乘积都等于 \(0^*\)。
\end{proposition}

\begin{proof}
加法的非严格包含直接由定义中的代表 \(a\in\alpha\subseteq\beta\) 得到。若 \(\alpha\subsetneq\beta\) 而平移后相等，即
\(\alpha+\gamma=\beta+\gamma\)，则两边再加 \(-\gamma\)，由结合律和逆元公式得到 \(\alpha=\beta\)，矛盾。因此加法严格保序。

下面设 \(\gamma>0^*\)；\(\gamma=0^*\) 的结论由吸收元公式得到。先固定 \(a_0\in\alpha\)，取有理数 \(t>-a_0+1\)。由于
\(t-\tfrac12\in t^*\) 且
\[
  a_0+t-\tfrac12>\tfrac12,
\]
分割 \(\alpha+t^*\) 含有正有理数。由于 \(\alpha\subseteq\beta\)，分割
\(\beta+t^*\) 也为正。正分割乘法的定义于是给出
\[
  (\alpha+t^*)\gamma\subseteq(\beta+t^*)\gamma.
\]
使用分配律展开，并在包含式两边加上 \(-(t^*\gamma)\)，得到
\[
  \alpha\gamma\subseteq\beta\gamma.
\]
若进一步有 \(\alpha<\beta\)，则由加法严格保序，
\[
  \beta-\alpha>0^*.
\]
两个正分割的积仍严格为正，因而
\[
 (\beta-\alpha)\gamma>0^*,
\]
再由分配律得到 \(\beta\gamma-\alpha\gamma>0^*\)，即
\(\alpha\gamma<\beta\gamma\)。这也说明为何严格保序必须排除
\(\gamma=0^*\)。
\end{proof}

\begin{theorem}
\((\mathbb D,+,\cdot,\subseteq)\) 是有序域。其零元、单位元分别为 \(0^*,1^*\)。
\end{theorem}

\begin{proof}
加法的封闭性、交换律、结合律、零元与逆元已在前文证明，所以 \((\mathbb D,+)\) 是交换群。正分割乘法演算给出正元素上的交换律、结合律和单位元；符号扩张把这些等式传递到任意符号，并且 \(0^*\) 为吸收元。一般分配律和每个非零元素的乘法逆元也已经证明。因此 \(\mathbb D\) 是域。

包含关系是全序。上一命题给出加法的保序性。若
\(0^*\subseteq\alpha\) 且 \(0^*\subseteq\beta\)，在该命题的乘法保序式中令较小元素为 \(0^*\)、乘数为 \(\beta\)，便有
\[
 0^*=0^*\beta\subseteq\alpha\beta.
\]
这正是序与乘法的相容条件，故 \(\mathbb D\) 是有序域。
\end{proof}

这里没有把关键的良定义问题隐藏在代表元之后：\(\mathbb D\) 的元素本身就是集合，所有运算均直接生成唯一集合。技术负担转移到了验证生成集合仍为分割以及逆元恒等式。

把全部域公理逐项列出可以看出，交换律和结合律主要继承自有理数运算；封闭性、加法逆元、乘法逆元和分配律需要分割条件与稠密性。尤其倒数构造必须使用右部信息，单靠正左部的逐项倒数会反转次序且不再向下封闭。

\section{确界性质}

至此已经得到一个有序域，但有序域仍可能像 \(\Q\) 那样带有缺口。分割构造的便利在这里显现：一族分割的最小上界无需另行逼近，它就是所有左部的并集。下面的证明也解释了为何最初选择“左部”作为新数的表示。


\begin{theorem}
有序域 \(\mathbb D\) 满足确界公理。
\end{theorem}

\begin{proof}
设 \(\mathcal A\subseteq\mathbb D\) 非空且有上界 \(\beta\in\mathbb D\)。令
\[
 \sigma=\bigcup_{\alpha\in\mathcal A}\alpha.
\]
非空性来自 \(\mathcal A\) 的任一成员；由于每个 \(\alpha\subseteq\beta\)，有 \(\sigma\subseteq\beta\ne\Q\)。并集保持向下封闭。若 \(x\in\sigma\)，则 \(x\in\alpha\) 对某个 \(\alpha\in\mathcal A\)；该 \(\alpha\) 无最大元，所以存在 \(y\in\alpha\subseteq\sigma\) 且 \(y>x\)。故 \(\sigma\) 是分割。

每个 \(\alpha\in\mathcal A\) 都包含于 \(\sigma\)，所以 \(\sigma\) 是上界。若 \(\gamma\) 是另一上界，则每个 \(\alpha\subseteq\gamma\)，从而其并 \(\sigma\subseteq\gamma\)。故 \(\sigma=\sup\mathcal A\)。
\end{proof}

\begin{corollary}[Dedekind 构造定理]
\(\mathbb D\) 是完备有序域，并且
\[
  \iota:\Q\longrightarrow\mathbb D,\qquad r\longmapsto r^*
\]
是保持 \(0,1\)、加法、乘法和次序的域嵌入。其像在
\(\mathbb D\) 中稠密。
\end{corollary}

\begin{proof}
有序域结构由前面各节的运算核验得到，确界定理由本节定理得到。主分割的加法、乘法及符号公式说明 \(\iota\) 保持域运算；保序性、单射性和稠密性已在第一节证明。
\end{proof}

这条推论完成了一个实数模型的构造。以后若采用分割模型，可以把
\(\iota(\Q)\) 直接记作 \(\Q\)，把 \(\mathbb D\) 记作 \(\R\)；这种记号上的识别以嵌入保持全部有序域结构为依据。分割的集合编码不会进入通常的极限计算，但确界的并集公式解释了这个模型为何没有 Dedekind 缺口。

\begin{proposition}
分割域中 \(1^*>0^*\)，每个非零分割的平方为正，并且不存在正分割小于所有正主分割 \(1/n^*\)。
\end{proposition}

\begin{proof}
前两项是任意有序域的形式后果。确界性质已经证明，因此可以应用\chapterref{chap:real-number-system}由确界原理推出的 Archimedes 性质：对任意正元素 \(\alpha\)，存在 \(n\in\Npos\) 使 \(1/n^*<\alpha\)。故不存在小于所有正主分割 \(1/n^*\) 的正分割。
\end{proof}

\section*{习题}
\addcontentsline{toc}{section}{习题}

\subsection*{A组　定义与基本方法}
\begin{enumerate}[label=\textbf{\thechapter.\arabic*.}]
 \exerciseitem{4.1} 判断下列集合是否为分割，并说明失败条件：
 \begin{enumerate}[label=(\arabic*)]
  \item \(\{q:q\le0\}\)；\item \(\{q:q<0\}\)；\item \(\Q\)；\item \(\{q:q<0\text{ 或 }q^2<2\}\)。
 \end{enumerate}
 \exerciseitem{4.2} 证明分割的右部非空且向上封闭，并刻画右部何时有最小元。
 \exerciseitem{4.3} 证明任意两个分割可比较，并说明证明中无最大元条件的作用。
 \exerciseitem{4.4} 证明 \(r\mapsto r^*\) 保持严格次序且为单射。
\end{enumerate}

\subsection*{B组　证明与反例}
\begin{enumerate}[label=\textbf{\thechapter.\arabic*.},resume]
 \exerciseitem{4.5} 完整验证 \(\alpha+\beta\) 的分割三条件。
 \exerciseitem{4.6} 对主分割直接计算 \(-r^*=(-r)^*\)。
 \exerciseitem{4.7} 对正有理数 \(r\) 直接计算 \((r^*)^{-1}=(r^{-1})^*\)。
 \exerciseitem{4.8} 证明若 \(\alpha\le\beta\)，则 \(\alpha+\gamma\le\beta+\gamma\)；若 \(0^*\le\gamma\)，则 \(\alpha\gamma\le\beta\gamma\)。
 \exerciseitem{4.9} 证明有界非空分割族的并是上确界，并解释交集为何一般不是下确界分割。
 \exerciseitem{4.10} 给出分割族，其交集有最大元；再写出修正后的下确界。
\end{enumerate}

\subsection*{C组　综合问题}
\begin{enumerate}[label=\textbf{\thechapter.\arabic*.},resume]
 \projectexerciseitem{4.11} \ 【前置：本章全部；预计用时：120分钟】设 \(F\) 是具有确界性质的有序域，并把 \(\Q\) 视为 \(F\) 中的有序子域。对 \(x\in F\) 定义
 \[
   L_x=\{q\in\Q:q<x\}.
 \]
 证明 \(L_x\) 是 Dedekind 分割，且映射
 \[
   \Psi:F\longrightarrow\mathbb D,\qquad x\longmapsto L_x
 \]
 是保持次序和域运算的双射。进而证明：在固定有理数嵌入以后，任意两个具有确界性质的有序域之间恰有一个保序域同构。
 \exerciseitem{4.12} 令 \(\alpha=\{q:q<0\text{ 或 }q^2<2\}\)。证明 \(\alpha^2=2^*\)。
 \optionalexerciseitem{4.13} 证明 \(\mathbb D\) 中任意非空有下界集合有下确界，并用上界或取负号给出两种表达。
 \openexerciseitem{4.14} 比较“分割没有最大元”和允许左部含边界的另一套约定；说明若不统一约定，有理数会出现两个表示。
 \exerciseitem{4.15} 证明分割加法严格保序；证明乘以正分割严格保序，并说明零因子边界。
 \exerciseitem{4.16} 补全主分割正数乘法
 \[
  r^*s^*=(rs)^*\qquad(r,s>0)
 \]
 的两个集合包含方向。
 \exerciseitem{4.17} 证明主分割在 \(\mathbb D\) 中稠密，即 \(\alpha<\beta\) 间存在 \(q^*\)。
 \exerciseitem{4.18} 证明若 \(\alpha<\beta\)，则存在正主分割 \(r^*>0^*\) 使 \(\alpha+r^*<\beta\)。
 \exerciseitem{4.19} 分别说明加法逆元与乘法逆元定义中为何使用右部元素，并构造一个“逐项取相反数或倒数”后不向下封闭的错误候选。
 \exerciseitem{4.20} 对非空有上界分割族 \(\mathcal A\)，验证并集上确界证明的每一条分割条件；再分析若删去共同上界假设哪一步失败。
 \optionalexerciseitem{4.21} 证明
 \[
  \inf\mathcal A=\bigcup\{\gamma\in\mathbb D:
  \gamma\subseteq\alpha\text{ 对所有 }\alpha\in\mathcal A\},
 \]
 并与“交集后必要时删端点”的公式比较。
 \projectexerciseitem{4.22} \ 【前置：本章全部运算；预计用时：120分钟】制作 Dedekind 构造的域公理核对表，逐项给出封闭性、单位元、逆元、结合律、分配律与序相容性的证明位置。
 \openexerciseitem{4.23} 比较以左部分割、右部分割和上下集合对三种方式构造实数的技术差异，说明三者怎样避免有理边界的双重表示。
 \optionalexerciseitem{4.24} 令
 \[
   \alpha=\{q\in\Q:q<0\text{ 或 }q^2<2\}.
 \]
 设 \((a_n)\) 是有理 Cauchy 列，并且从某项起 \(a_n>0\)。再假设对每个正有理数 \(\varepsilon\)，存在 \(N\) 使
 \[
   n\ge N\quad\Longrightarrow\quad \abs{a_n^2-2}<\varepsilon.
 \]
 定义
 \[
 \begin{aligned}
   L_{(a_n)}&=\{q\in\Q:\exists N\ \forall n\ge N,\ q<a_n\},\\
   R_{(a_n)}&=\{q\in\Q:\exists N\ \forall n\ge N,\ a_n<q\}.
 \end{aligned}
 \]
 证明 \(L_{(a_n)}=\alpha\) 且 \(R_{(a_n)}=\Q\setminus\alpha\)。因此，任意满足上述条件的有理列都给出同一个 Dedekind 分割。
\end{enumerate}

\section*{部分习题提示}
\begin{itemize}
 \item \exerciseref{4.7}：把 \(q<1/r\) 与存在 \(s\ge r\) 使 \(q<1/s\) 对照。
 \item \exerciseref{4.10}：若交集有最大元 \(c\)，应删除 \(c\) 以恢复分割条件。
 \item \exerciseref{4.12}：先证每个正乘积 \(ab<2\)，再从 \(x<2\) 构造足够接近缺口的 \(a,b\) 使 \(x<ab\)。
\end{itemize}
